\subsection{Bounded negative action and the central source interval}
\label{subsec:bounded-action}

The bounded-footpoint passage and the central source passage use different
sections and phase labels. We connect them by following the same orbits
between those sections. This gives the overlap needed to combine a
negative-curvature estimate with a bounded-derivative estimate.

Let \(\tau>0\), \(\sigma>0\), and
\begin{equation}\label{bd:polynomial}
\begin{split}
 P_\sigma(\xi)
 &=(1+\sigma\xi)(1-\xi/x_*)(1-\sigma j\xi)
   =1-\Lambda_1\xi+3\Lambda_3\xi^2+\Lambda_2\xi^3,\\
 \Lambda_1&=x_*^{-1}-\sigma+\sigma j,\qquad
 \Lambda_2=\sigma^2j/x_*,\qquad
 3\Lambda_3=\sigma(j-1)/x_*-\sigma^2j .
\end{split}
\end{equation}
Here \(x_*\) ranges in a positive compact interval and \(j\) in a
fixed compact subset of \([0,\infty)\).  Physical raw sections require
\(j>0\).  The value \(j=0\) denotes only a one-sided normalized limit.
We take either a positive compact \(j\)-interval, or a sufficiently
small interval \(0\le j\le j_0<1\).  In the latter case
\(\Lambda_3<0\).  In both cases, after decreasing \(\sigma\),
\[
 \Lambda_1\ge\lambda_1>0,\qquad
 0\le\Lambda_2\le C\sigma^2,\qquad |\Lambda_3|\le C\sigma.
\]
No comparison between \(\tau\) and \(\sigma\), or between either of
them and \(j\), is imposed.

The source parameters and physical coordinates are
\begin{equation}\label{bd:source-dictionary}
\begin{gathered}
 (a,p_{\rm s},\gamma_{\rm s},s_{\rm s})
 =(-1,\sqrt\tau\Lambda_1,\tau\Lambda_3,\tau^{3/2}\Lambda_2),\\
 \xi=\sqrt\tau X,\quad\eta=\tau Y,\qquad
 v-\gamma_{\rm s}=e\xi,\quad z=e^2\eta,\quad e=\tau\Lambda_2.
\end{gathered}
\end{equation}
For a source quartet, denote its positive graph denominator by
\[
 \mathcal D_{a,p,\gamma,s}(X,Y)
 =\{1-A(\gamma+sX)+C(\gamma+sX)^2\}Y
       -a-pX+3\gamma X^2+sX^3.
\]
We write \(\mathcal D_{\rm s}\) for the quartet in
\eqref{bd:source-dictionary}; barred subscripts below mean the same
formula with the displayed barred quartet. These denominators determine
where an orbit can be written as a graph.
Writing
\[
 B_\tau=1+\tau(2\Lambda_3+\Lambda_2\xi),\quad
 L_\tau=\Lambda_3+\Lambda_2\xi,\quad
 Q_\tau=1-A\tau L_\tau+C\tau^2L_\tau^2,
\]
the exact graph equation is
\begin{equation}\label{bd:graph}
 \eta_\xi=
 \frac{\eta\{\xi B_\tau+C\tau^2\Lambda_2L_\tau\eta\}}
      {Q_\tau\eta+\tau P_\sigma}.
\end{equation}
Fix \(L>4(1+\sup x_*)\).  On the actual incoming section \(\xi=L\)
use the increasing coordinates
\begin{equation}\label{bd:incoming-phase}
 h_L=\eta+\tau\log(2\eta)-L^2/2,
 \qquad \beta_L=h_L/\tau-\log(2\tau).
\end{equation}
These formulas prescribe initial points; neither expression is asserted
to be constant along a perturbed orbit.  In particular,
\(h_L=0\), or \(\beta_L=b_\tau:=-\log(2\tau)\), is an actual
interior source orbit.

\begin{lemma}[Source ranges on the actual incoming section]
\label{bd:lem-source-ranges}
There is a fixed \(\delta_\sharp>0\) such that the following
constructions are uniform on \(-\delta_\sharp\le h_L\le0\), with
enlarged input and output intervals.
\begin{enumerate}
\item For \(B_0\le\beta_L\le b_\tau\), where \(B_0\) is fixed
large, the source passage exists and its actual central height obeys
\begin{equation}\label{bd:positive-height}
 |Y(0)+\log Y(0)-\beta_L|\le C_L,
 \qquad Y(0)=\beta_L-\log\beta_L+O_L(1).
\end{equation}
\item Every fixed compact \(\beta_L\)-interval reaches a fixed
compact strict source family: on fixed internal source intervals,
\begin{equation}\label{bd:compact-height}
 0<m_B\le Y(0)\le M_B,\qquad
 Y>0,\qquad \mathcal D_{\rm s}\ge c_B(Y+1)>0.
\end{equation}
\item At a negative base \(\beta_L=-B\), \(B\ge B_-\), freeze
\begin{equation}\label{bd:frozen-dictionary}
 q=B^{-1/2},\quad \rho=\sqrt{\tau B},\quad
 (x,y)=(\xi/\rho,\eta/\rho^2),\quad
 (\lambda_1,\lambda_2,k)
       =(\rho\Lambda_1,\rho^3\Lambda_2,\rho^2\Lambda_3).
\end{equation}
For fixed cuts \(x=\pm L_{\rm f}\), \(L_{\rm f}=8\), this orbit
belongs to a common enlarged negative-action family with core
\(|x|\le R_{\rm f}=5\) and \(P_\sigma(\rho x)\ge3/4\).
The input and output energy windows can be taken respectively inside
\([-7/4,-3/8]\) and \([-2,-1/4]\), with strict margins.
\end{enumerate}
All these same-orbit passages reach \(\xi=\pm L\) with
\begin{equation}\label{bd:parabolic-cuts}
 \eta(\pm L)=L^2/2+
 O_L\bigl(\delta_\sharp+\tau|\log\tau|\bigr).
\end{equation}
Their outward continuations to the fixed physical sections have
positive graph denominators and retain the actual regular maps.
\end{lemma}
\begin{proof}
For the first assertion, freeze \(B=\beta_L\ge B_0\), put
\(q=B^{-1/2}\), \(\rho=\sqrt{\tau B}\), and use the source
scaling \((\xi,\eta)=(\rho x,\rho^2y)\).  Its coefficients are
\[
 (\bar a,\bar p,\bar\gamma,\bar s)
 =(-q^2,\sqrt\tau\Lambda_1q,\tau\Lambda_3,
                         \tau^{3/2}\Lambda_2/q).
\]
Set \(H=y+q^2\log(2y)-x^2/2\) and \(R=L/\rho\).
On a provisional quadratic corridor \(y\asymp1+x^2\), direct
subtraction in the exact graph equation gives
\begin{equation}\label{bd:energy-estimates}
\begin{split}
 |H_x|&\le C_L\{|
 \bar p|+|\bar\gamma|(1+|x|)+\bar s(1+x^2)\},\\
 |\partial_HH_x|&\le C_L\left\{
 \frac{|\bar p|}{1+x^2}+
 \frac{|\bar\gamma|}{1+|x|}+\bar s\right\}.
\end{split}
\end{equation}
The terms with \(\bar p=\bar\gamma=\bar s=0\) cancel exactly.
The first integral on \([-R,R]\) is \(O_L(q^2)\), since its three
largest products are
\(L|\Lambda_1|q^2,L^2|\Lambda_3|q^2,L^3\Lambda_2q^2\).
The second integral tends to zero uniformly for
\(B_0\le B\le b_\tau\).  Start at \(x=0\) with
\(H(0,y)\in[1/4,3]\).  The monotonicity of
\(y+q^2\log(2y)\) improves the bounds at any first contact with the corridor boundary:
\[
 (1+x^2)/16<y<4(1+x^2),\qquad \mathcal D_{\rm bar}\ge cy.
\]
Thus both cuts are reached.  The scalar energy variation is a positive
exponential with exponent bounded by the second integral in
\eqref{bd:energy-estimates}; hence the incoming energy map has an
actual inverse on a common enlarged interval.  Its prescribed value at
\(x=R\) is exactly \(1+B^{-1}\log(2/B)\).  Inverting and undoing
the scale gives \eqref{bd:positive-height}.  This also constructs the
orbit \(h_L=0\) and a real larger source chord through it.

For compact \(\beta_L\), work instead in \((X,Y)\) and put
\(H_0=Y+\log(2Y)-X^2/2\), \(R_0=L/\sqrt\tau\).
At the incoming cut, \(H_0=\beta_L+\log2\).  On a provisional
bounded energy region the denominator is at least \(c(Y+1)\): near
\(\xi=0\), \(P_\sigma\ge1/2\), while away from that fixed small
interval bounded energy forces \(Y\ge X^2/4\), which dominates
the bounded polynomial.  The exact energy subtraction gives
\[
 |H_0'|\le C_L\{p_{\rm s}+|\gamma_{\rm s}|(1+|X|)
                +s_{\rm s}(1+X^2)\}(1+|H_0|).
\]
Its coefficient has a bounded integral because
\(p_{\rm s}R_0=L\Lambda_1\),
\(|\gamma_{\rm s}|R_0^2=L^2|\Lambda_3|\), and
\(s_{\rm s}R_0^3=L^3\Lambda_2\).  Gronwall's inequality improves a
fixed sufficiently large energy bound, and proves continuation.
Inverting \(Y+\log(2Y)\) at the centre proves
\eqref{bd:compact-height}.  Choose larger compact incoming labels
first to obtain strict image margins.  The resulting internal energy
coordinate may retain a bounded translation; no identity with
\(\beta_L+\log2\) on another section is used.

For the negative base in \eqref{bd:frozen-dictionary}, the energy
\(H_{\rm f}=y+q^2\log(2y)-x^2/2\) at \(x=L/\rho\) is
\begin{equation}\label{bd:frozen-incoming}
 b_B(\beta)=q^2\beta+q^2\log(2q^2),\qquad
 b_B(-B)=-1-\frac{\log(B/2)}B.
\end{equation}
Also \(\rho^2=-h_L+\tau\log(2\tau)\le-h_L\).
Choose \(\delta_\sharp\) so that
\(P_{\rm f}(x):=P_\sigma(\rho x)\ge3/4\) on
\([-8,8]\).  On the positive quadratic tail the exact subtraction
gives
\[
 |(H_{\rm f})_x|
 \le Cq^2\{\rho\Lambda_1+
             \rho^2|\Lambda_3||x|+\rho^3\Lambda_2x^2\}.
\]
Its integral to \(L/\rho\) is \(O_L(q^2)\).  A first-contact
argument in a quadratic corridor therefore reaches \(x=8\), with
the actual energy in the interior of \([-3/2,-1/2]\), after increasing \(B_-\). This places the orbit in the negative-action
family. Its relative curvature estimate will be proved below.

To obtain phase derivatives on a neighbourhood of these orbits, we
construct the full negative core on the larger input window
\([-7/4,-3/8]\).  Fix the finite derivative budget \(m\) and an
integer \(M>2m+8\).  Follow the positive fast arc to
\(y=q^{M+2}\).  The logarithmic-height equation has
\(dH_{\rm f}/d\log y=O_{C^m}(q^2\rho)\); its clock length is
\(O(|\log q|)\).  It cannot reach \(x=1/2\) first, since that
would imply \(H_{\rm f}\ge-1/8+o(1)\).  Hence its actual lip
coordinate satisfies
\[
 p=\sqrt{-2b}+O_{C^m}(q^2|\log q|),\qquad 1/2<p<2.
\]
In the thin core put \(\mathfrak q=q^2\log(y/q^2)\) and
\[
 Q_{\rm f}(x)=1-Aq^2(k+\lambda_2x)
                        +Cq^4(k+\lambda_2x)^2.
\]
Until its first return to the same lip height, its exact equation is
\begin{equation}\label{bd:frozen-core}
 \mathfrak q_x=
 \frac{x\{1+q^2(2k+\lambda_2x)\}
       +Cq^6\lambda_2(k+\lambda_2x)e^{\mathfrak q/q^2}}
      {P_{\rm f}(x)+Q_{\rm f}(x)e^{\mathfrak q/q^2}}
 =\frac{x}{P_{\rm f}(x)}+O(q^2\rho)+O(q^M).
\end{equation}
The actions \(\int_0^p v/P_{\rm f}(v)\,dv\) and
\(\int_0^r v/P_{\rm f}(-v)\,dv\) have opposite comparison
differences at \(r=p/2\) and \(r=2p\), with fixed margins.
The logarithmic height first decreases, stays a fixed amount below
the lip near zero, then increases strictly on the negative half.
It therefore has one first exit before \(|x|=5\), with
\[
 r=p+O(\rho+q^2|\log q|)>1/4.
\]
The denominator in \eqref{bd:frozen-core} is at least \(3/4\).
Each exponential derivative costs at most two powers of \(q^{-1}\);
the choice of \(M\) and the transverse exit equation give bounded
finite phase derivatives on nested windows.  The negative fast arc
then reaches \(x=-8\), with
\(b_-=b+O_{C^m}(\rho+q^2|\log q|)\) strictly inside
\([-2,-1/4]\).  The exact independent energy coordinates of
Theorem~\ref{neg:thm-relative} differ from \(H_{\rm f}\) by
\(O_{C^m}(|k|)\) at these fixed positive-height cuts.  Slightly larger
windows thus contain those coordinates and their images as well.

The same energy estimates on both outward tails prove
\eqref{bd:parabolic-cuts}.  In raw variables, the quadratic corridor
of \eqref{eq:parabolic-tail} has a denominator comparable to the
squared radial variable.  Choose \(L\) large and
\(\delta_\sharp\) small first.  All the enlarged chords above enter
its interior, so the first-contact argument continues them to the
fixed raw sections, across any intervening axis root at positive
height.  The fixed regular transitions are then included.  These
are actual first hits of the original field throughout.
\end{proof}

\begin{lemma}[The bounded annulus]
\label{bd:lem-annulus}
Fix positive constants \(p_-,r_{\rm a},R_{\rm BD}\), with
\(p_-<\inf x_*\).  Consider one connected family with positive and
negative strong-axis footpoints
\[
 p_-\le p<x_*,\qquad r_{\rm a}/2\le r\le2R_{\rm BD},
 \qquad b_+=-p^2/2,\quad b_-=-r^2/2.
\]
The projections are synchronized to fixed weak sections, and their
outward continuations use the fixed physical sections of
Lemma~\ref{lem:closing}.  On slightly larger such domains, normalize
each actual embedding by its own exact value and positive first
derivative.  The resulting complete map
\(\Pi=\phi_-\circ S\circ\phi_+^{-1}\) satisfies
\begin{equation}\label{bd:annular-curvature}
 \phi_+'\frac{P_\sigma(p)}{p+r}\frac{\Pi''}{\Pi'}
 =3\Lambda_2-\frac{\Lambda_1}{pr}+E,\qquad
 \|E\|_{C^1_{b_+}}\longrightarrow0 .
\end{equation}
The derivative of the error follows the actual matched family.
That family has a uniform positive gap from the positive root, as
proved below; convergence includes the reciprocal coefficient face
\(j=0\). For sufficiently small coefficients this map is strictly concave and its displacement
has at most two zeros.
\end{lemma}
\begin{proof}
The curvature formula is uniform only after the incoming footpoint has
been separated from the positive root. We first obtain this separation
by matching the two normal heights. On fixed
small strong lips the normal rates are positive units, including at
\(p=x_*\), since
\[
 -P_\sigma'(x_*)\ge c_*,\qquad
 P_\sigma(-r)\ge c_-,\qquad rB_\tau(-r)\ge c_- .
\]
The strong projection and its integrating normal coordinate are uniform
through the required finite phase order: the backward tangential Green
factor is at most \(Ce^{C\tau t}\), whereas the normal variations are
\(C(1+t)^Ne^{-ct}\). Their integral therefore defines the projection
and the coordinate satisfying exactly \(\dot z=-pB_\tau(p)z\) on
the positive side, with the corresponding equation in reversed time on the negative side.
Use the original fixed weak cuts \(c_+<p_-\) and \(c_-<r_{\rm a}/2\).
Integration along the projected axis gives their actual normal heights
\[
 w_+(p)=u_+(p)e^{-J_+(p)/\tau},\qquad
 w_-(r)=u_-(r)e^{-J_-(r)/\tau},\qquad u_\pm^{\pm1}\le C,
\]
where the units and their finite phase derivatives are bounded on the
fixed lips, independently of any outer cut, and
\[
 J_+(p)=\int_{c_+}^p\frac{xB_\tau(x)}{P_\sigma(x)}\,dx,
 \qquad J_-(r)=\int_{c_-}^r\frac{xB_\tau(-x)}{P_\sigma(-x)}\,dx .
\]

On the intervening weak interval \(P_\sigma,Q_\tau\ge c>0\), so
\(-\dot\xi=\tau P_\sigma+Q_\tau\eta\ge c\tau\) and its travel
time is at most \(C/\tau\). Put
\(H=\xi B_\tau+C\tau^2\Lambda_2L_\tau\eta\), so
\(d\log\eta/dt=-H\). At a critical height,
\(d^2\log\eta/dt^2=H_\xi(\tau P_\sigma+Q_\tau\eta)>0\).
The standing passive neighborhood has \(C>0\), and
\(H_\xi=1+\tau(2\Lambda_3+2\Lambda_2\xi)
+C\tau^2\Lambda_2^2\eta>1/2\) on this word. Thus there is no
interior height maximum; the buffered weak endpoint heights bound the
whole passage and its logarithmic rate. At each fixed weak section
the integrating normal coordinate and \(\eta\) differ by a positive
unit. For the actual map \(M_\tau\) between those normal coordinates,
\[
 \left|\tau\log\frac{M_\tau(w)}w\right|\le K .
\]
Its positive first multiplier \(M_\tau'(0)\) is retained. Matching
the two heights now bounds \(J_+(p)\) by \(\sup J_-(r)+K+C\tau\).
The latter is uniform for \(r\le2R_{\rm BD}\), whereas the simple
positive root gives
\[
 J_+(p)\ge a_*\log\frac{\ell_*}{x_*-p}-C,
 \qquad a_*>0 .
\]
The logarithmic residue \(x_*B_\tau(x_*)/[-P_\sigma'(x_*)]\)
is bounded below by \(a_*\).
Consequently every actual member has \(p\le x_*-d_{\rm a}\), for
one fixed \(d_{\rm a}>0\). This gap depends only on the fixed lips and middle passage, so the
outer cuts can now be chosen without changing it.

Enlarge the input within \(p<x_*-d_{\rm a}/2\). The limiting
root-free action has a fixed compact output envelope: at \(\sigma=0\),
its negative action grows linearly at infinity. Choose strict lower
and upper exit brackets there, then decrease \(\sigma\) independently
so they and their polynomial margins persist. Denote an enlarged
upper output bound by \(R_{\rm a}\ge2R_{\rm BD}\). We can therefore choose the annular cuts
\[
 L_{\rm a}>4(1+\sup x_*+R_{\rm a}).
\]
The source cut \(L\), its label \(h_L\), and
\(\delta_\sharp(L)\) stay fixed. The positive-height arcs from the
fixed lips to \(\pm L_{\rm a}\) and their small enlarged chords
have uniform height and denominator margins. Independent final
coefficient bounds place \(eL_{\rm a}\) inside the original raw
ports; scalar uniqueness identifies both decompositions there.

Apply the compact coefficient-annulus argument in the proof of
Corollary~\ref{neg:cor-rootfree}, with \(t^2=\tau\),
\(\lambda=(\Lambda_1,\Lambda_2)\), and \(k=\Lambda_3\).
The root-free polynomial margin, strict exit brackets and nonzero
negative exit slope established above are its inputs. The lip and
first-exit estimates give absolute \(C^3\) error
\(O(\tau(1+|\log\tau|))\), including synchronized resection;
no division by the reversible coefficient distance is needed. If
\(S_0(-p^2/2)=-r_0^2/2\) is defined by the exact limiting action,
\begin{equation}\label{bd:annular-first}
 \|S-S_0\|_{C^3}\longrightarrow0,\qquad
 S_0'=\frac{P_\sigma(-r_0)}{P_\sigma(p)} .
\end{equation}
The exact polynomial identity therefore gives
\[
 \frac{P_\sigma(p)}{p+r_0}\frac{S_0''}{S_0'}
 =\frac{P_\sigma'(p)/p+P_\sigma'(-r_0)/r_0}{p+r_0}
 =3\Lambda_2-\frac{\Lambda_1}{pr_0} .
\]

The two outer embeddings may have very different first derivatives.
We estimate each embedding relative to its own exact first derivative.
On a fixed small positive-height lip let \(p(t)\) be the synchronized
axis solution and \(w=\eta-\xi^2/2+p(t)^2/2\).
The exact field gives
\[
 \dot w=\tau\{\xi P_\sigma(\xi)-pP_\sigma(p)\}
 +\xi\eta(Q_\tau-B_\tau)
 -C\tau^2\Lambda_2L_\tau\eta^2.
\]
Since \(|\xi+p|\) is bounded below, this has the form
\(\dot w=\tau a(t)w+\tau\eta d(t)\), with bounded finite
derivatives and exponentially decreasing \(\eta\).  The fixed
weak section is a positive axis distance away and the axis speed is
\(O(\tau)\), so its remaining time is at least \(c/\tau\).
Terminal hitting-time derivatives multiply the same exponentially
small normal height.  Variation of constants, with four phase
derivatives, gives \(w=O_{C^4}(\tau)\); resecting a moving lip
costs \(O_{C^4}(\tau|\log\tau|)\).  On both fixed outer cuts,
therefore,
\begin{equation}\label{bd:annular-tail-data}
 \eta_\pm(\pm L_{\rm a},b)=L_{\rm a}^2/2+b+
                  O_{C_b^4}(\alpha_\tau),\qquad
 \alpha_\tau=\tau(1+|\log\tau|).
\end{equation}

Set \(s=e\xi\), \(Z=e^2\eta\), \(k_0=\Lambda_3\).
The exact raw tail equation is
\[
 Z_s=\frac{n_2Z^2+n_1Z}{\ell_oZ+m_o},\quad
 \begin{aligned}
 \ell_o&=1-A(\tau k_0+s)+C(\tau k_0+s)^2,\\
 m_o&=s^3+3\tau k_0s^2-\tau\Lambda_1es+\tau e^2,\\
 n_1&=s(1+s+2\tau k_0),\quad n_2=C(\tau k_0+s).
 \end{aligned}
\]
The initial data \eqref{bd:annular-tail-data} enter the invariant
corridor \(Z\asymp s^2\), \(\ell_oZ+m_o\asymp s^2\).
Division writes the right side as
\(F=(n_2/\ell_o)Z+q_o+\mathfrak r/(\ell_oZ+m_o)\), where
\[
 |\mathfrak r|\le C(|s|^4+\tau|k_0||s|^3
                       +\tau e|s|^2+\tau e^2|s|).
\]
If \(j\) is bounded away from zero,
\(\tau|k_0||\log e|\to0\), because
\(k_0=O(\sigma)\) and \(e\asymp\tau\sigma^2\).
Thus the first variation \(J_1=Z_b\) is comparable to \(e^2\).
For the reciprocal range \(j\to0\), retain instead \(k_0<0\).
Direct differentiation gives the exact first cocycle
\[
 F_Z=\frac{n_2\ell_oZ^2+m_o(n_1+2n_2Z)}{(\ell_oZ+m_o)^2},
 \qquad \operatorname{sgn}(s)(n_1+2n_2Z)\ge c|s| .
\]
The latter follows from \(Z\asymp s^2\) and the small fixed raw
radius. At the actual common denominator,
the term \(3\tau k_0s^2\) in \(m_o\) contributes nonpositively
to \(\operatorname{sgn}(s)F_Z\) on both outward tails.  The rest gives
\[
 \operatorname{sgn}(s)F_Z\le
 C(1+\tau e/|s|^2+\tau e^2/|s|^3),\qquad 0<J_1\le Ce^2.
\]
The normalization below uses this exact positive derivative and requires
only the displayed upper bound. For \(2\le n\le3\), in both cases,
\[
 \int |F_{Z^n}|J_1^{n-1}\,d|s|
 \le C(eL_{\rm a}^{3-2n}+\tau|k_0|L_{\rm a}^{2-2n}
                       +\tau L_{\rm a}^{1-2n}+\tau L_{\rm a}^{-2n})\le C\tau.
\]
The quantitative normalization following
\eqref{tail:relative-jet-equations}, through order three, now applies
with initial errors from \eqref{bd:annular-tail-data} and
\(\int|F_{Z^n}|J_1^{n-1}=O(\tau)\), \(n=2,3\).
Each tail and its actual-image inverse are consequently
\(\mathrm{id}+O_{C^3}(\alpha_\tau)\) after normalization by
the retained \(J_1(b_*)>0\). At this actual image the ordinary
\(C^3\) observer adds \(O_{C^3}(J_1(b_*))\) by
\eqref{phys:eq-observer-normalization}. Thus
\(\phi_\pm=\mathrm{id}+o_{C^3}(1)\), including when their
individual gains are much smaller than \(e^2\).

Finally the exact chain rule is
\[
 \frac{\Pi''}{\Pi'}=\frac1{\phi_+'}
 \left\{\frac{S''}{S'}+
       \frac{\phi_-''(S)}{\phi_-'(S)}S'
       -\frac{\phi_+''}{\phi_+'}\right\}.
\]
On the compact family, \(S\), its first three derivatives and
its positive first inverse derivative are bounded. Multiply by
\(P_\sigma(p)/(p+r)\) and use \(r-r_0\to0\) in phase
\(C^3\). Both exterior corrections are \(o_{C^1_{b_+}}(1)\).
This proves \eqref{bd:annular-curvature}.  Its principal term is at most
\(3C\sigma^2-\lambda_1/(2(\sup x_*)R_{\rm BD})<0\).
Positive affine changes preserve this sign in the actual input.
Thus \(x-m\Pi(x)\) is strictly convex, proving the two-zero bound.
\end{proof}

\begin{proposition}[Bounded-footpoint domains and derivatives]
\label{bd:prop-bounded-action}
On one ordered through interval ending at the actual cut \(h_L=0\),
suppose the two fixed comparisons described below and the enlarged
annular domains of Lemma~\ref{bd:lem-annulus} use the same actual
physical sections. The union of the central source interval
\(-\delta_2\le h_L\le0\) and the part with negative strong
footpoint \(0<r\le2R_{\rm BD}\) has an overlapping cover by the
bounded annulus and the negative, compact, and positive source-label
intervals constructed below. All endpoints are fixed before choosing
a displacement zero.

The completed map has strictly negative curvature on the annulus and
on the entire negative source interval. On the compact and positive
source intervals its derivative in the actual affine input is bounded
by one constant. After one common multiplier shrink, the displacement
has positive derivative on these latter intervals, including the cut
\(h_L=0\). The same assertions hold on the enlarged intervals used
for the overlaps. The degree-loss passage below verifies all domain
and overlap hypotheses of this proposition.
\end{proposition}
\begin{proof}
Fix the strict source domain and the constants for its finite phase
derivatives. Choose $B_-$ large, then $\delta_\sharp$ small enough
for the relative negative-curvature estimate below. Choose a fixed $0<\delta_1\le\delta_\sharp/8$,
set $p_0=\sqrt{2\delta_1}$, and require $P_\sigma\ge1/2$ on
$[-2p_0,2p_0]$.
The orbit prescribed by \(h_L=-\delta_1\) is constructed by a
compact fast arc to a small lip, followed by the core first-exit
argument in \eqref{bd:frozen-core}.  Its limiting incoming footpoint
is \(p_0\), and the limiting outgoing footpoint is determined by
\[
 \int_0^r\frac{v}{P_\sigma(-v)}\,dv
       =\int_0^{p_0}\frac{v}{P_\sigma(v)}\,dv.
\]
The comparison values \(r=p_0/2,2p_0\) have strict opposite action
differences.  Thus the actual two footpoints lie in
\((p_0/2,2p_0)\) after shrinking \(\tau\).
Choose \(r_{\rm a}<p_0/4\).  A second fixed comparison
\(h_L=-p_{\rm a}^2/2\), with \(0<p_{\rm a}<r_{\rm a}/8\), has
\(r<r_{\rm a}/4\) and \(p>p_{\rm a}/2\).
Scalar order and the actual synchronized projections imply that
every \(r\ge r_{\rm a}/2\) in the bounded annulus has
\(p\ge p_{\rm a}/2\), while every \(r<r_{\rm a}\) lies in
\(-\delta_1<h_L\le0\).  Both comparisons have strict margins,
so these ranges have an open overlap.  The comparison uses order along the same orbits; \(h_L\) need not equal
the footpoint energy on the other section.

Fix \(\delta_1<\delta_2<\delta_\sharp\).  On
\(-\delta_2\le h_L\le0\), divide the actual incoming labels into
\[
 [\beta_{\min},-B_-],\quad[-B_-,B_+],\quad[B_+,b_\tau],
 \qquad \beta_{\min}=-\delta_2/\tau-\log(2\tau).
\]
Each interval is used with slightly larger fixed phase thresholds.
For a negative base, Lemma~\ref{bd:lem-source-ranges} supplies the
independent enlarged windows $\widehat J_+=[-7/4,-3/8]$ and
$\widehat J_-=[-2,-1/4]$, frozen cuts $L_f=8$ and core $R_f=5$.
The actual incoming frozen energy is strictly inside $[-3/2,-1/2]$;
its separate ideal coordinate differs by $O(|k|)$. Choose a fixed
$J\Subset\widehat J_+$ containing these values with a strict buffer.
With the frozen dictionary~\eqref{bd:frozen-dictionary}, put
\[
 \rho_{\rm id}=(\lambda_1^2+\lambda_2^2)^{1/2}.
\]
Use the tail part of the proof of Theorem~\ref{neg:thm-relative}
with \(t=q\), \(\lambda=(\lambda_1,\lambda_2)\) and the frozen
\(k\). Its coefficients are exactly
\((a,p,\gamma,s)=(-q^2,q^2\lambda_1,q^2k,q^2\lambda_2)\),
and its auxiliary cuts are \(v=q^2k\pm d_H\), with fixed \(d_H\).
The enlarged windows above supply its quadratic corridors;
at \(L_f=8\) its fourth-order kernel majorants become small after
a \(q\) shrink. The same sequential intrinsic \(p\)- and
\(s\)-secants, with the two reference labels fixed, give
\begin{equation}\label{bd:intrinsic-relative-tails}
 \|\phi_\pm^{\rm int}-\mathrm{id}\|_{C^3}
 +\|(\phi_\pm^{\rm int})^{-1}-\mathrm{id}\|_{C^3}
 \le Cq^2\rho_{\rm id}.
\end{equation}
For the actual ordinary \(C^3\) observers, including divided closing,
the separate raw gains satisfy, for fixed \(0<\nu<1\),
\begin{equation}\label{bd:relative-observer-gain}
 0<d\le Cs_f^{2-\nu}\le Cs_f
   =Cq^2\lambda_2\le Cq^2\rho_{\rm id},\qquad
 s_f=q^2\lambda_2=e\rho .
\end{equation}
Equation~\eqref{phys:eq-observer-normalization} on the fixed
buffered cut-label chord therefore preserves
\eqref{bd:intrinsic-relative-tails} for the observed tails and their
actual-image inverses. The fourth phase order is used only in the
intrinsic polynomial kernel, before composing with these observers.

On the negative source interval the limiting curvature tends to zero
with the reversible coefficient distance. We therefore compare the
actual core \(S_q\) with its limiting action \(S_0\) before dividing
by that distance. The exact reversible energy gives
\(S_q(0,k)=S_0(0,k)=\mathrm{id}\). Integrating the coefficient
variation in \eqref{neg:compact-mixed-convergence} before taking any
limit yields
\[
 \|S_q-S_0\|_{C^2}
 \le C\rho_{\rm id}q^2(1+|\log q|).
\]
One intrinsic coefficient secant with two phase derivatives produces
the factor \(\rho_{\rm id}\). This makes the error small relative
to the core curvature even as that curvature vanishes.
For \(b\in J\), put \(p=\sqrt{-2b}\), and let \(r_0\) be its
limiting action exit. Equation~\eqref{neg:action-curvature} gives
\[
 \frac{S_0''(b)}{S_0'(b)}
 =\frac{p+r_0}{P_{\rm f}(p)}
       \left(3\lambda_2-\frac{\lambda_1}{pr_0}\right).
\]
The even coefficient \(k\) cancels exactly. Here
\(\rho^2=\tau B\le-h_L\), \(\rho_{\rm id}\asymp\rho\),
\(\lambda_1/\rho_{\rm id}\ge c_1>0\), and
\(\lambda_2/\rho_{\rm id}=O(\rho^2\sigma^2)\).
The compact root-free margins and positive first variation thus give
\(S_0''/\rho_{\rm id}\le-c_0<0\), without a Taylor error in \(k\).

The complete curvature also contains the two section terms:
\[
 \Pi''(\phi_+(b))=
 \frac{\phi_-'(S_q(b))}{\phi_+'(b)^2}
 \left\{S_q''+(S_q')^2\frac{\phi_-''(S_q)}{\phi_-'(S_q)}
                  -S_q'\frac{\phi_+''}{\phi_+'}\right\}.
\]
The separate tail bounds above make the two section terms
\(O(q^2\rho_{\rm id})\), while the prefactor is a positive unit
and differs from one by that amount. Consequently, at the actual
normalized input \(u_*=\phi_+(b)\),
\begin{equation}\label{bd:relative-negative-curvature}
 \frac{\Pi''(u_*)}{\rho_{\rm id}}
 =\frac{S_0''(b)}{\rho_{\rm id}}
       +O\bigl(q^2(1+|\log q|)\bigr)<-c .
\end{equation}
Choose \(B_-\) large enough to absorb this error, then
\(\delta_\sharp\) small enough for the root-free domain and the
fixed relative coefficient margins. No comparison of \(q\) with \(\rho\) is used.
The incoming resection identifies the point $u_*$ at which this sign
is evaluated. For the two retained positive first gains,
\begin{equation}\label{bd:relative-raw-curvature}
 (\Pi^{\rm raw})''=\frac{g_-}{g_+^2}\Pi'',\qquad g_+,g_->0.
\end{equation}
Each negative base is fixed while differentiating, so this is one
strict sign on the whole negative interval, including nonroot inputs.

On the compact middle interval, \eqref{bd:compact-height} and the
larger label windows give a compact strict core and positive variation
units $A_\pm$. Before normalization its two auxiliary raw variations are
\begin{equation}\label{bd:compact-raw-variation}
 \frac{dz_\pm}{dY_0}
 =s_s^{2-\lambda_\gamma}(d_{H,\pm}/L_t)^{\lambda_\gamma}A_\pm U_\pm,
 \quad s_s=\tau^{3/2}\Lambda_2,\quad
 \lambda_\gamma=\frac{12\gamma_s}{1+2\gamma_s},\quad
 A_\pm^{\pm1},U_\pm^{\pm1}\le C.
\end{equation}
The common singular power cancels in $dz_-/dz_+$; distinct fixed
port radii have a bounded ratio. The actual regular bridges and divided
closing have positive derivative units, so $0<\Pi'\le C_B$.
This cancellation gives the derivative bound without a lower bound on a
raw gain or a bound on $\gamma_s\log s_s$.
On the positive interval, \eqref{bd:positive-height} gives, after
choosing \(B_+\) large,
\[
 Y(0)>K(1+p_{\rm s}^2),\qquad
 0<s_{\rm s}^2Y(0)\le C_L\tau^3\sigma^4j^2|\log\tau|<z_*.
\]
Since \(\alpha_0=-s_{\rm s}^2\) and
\(\alpha_1=s_{\rm s}p_{\rm s}\), these are the hypotheses of
Corollary~\ref{rew:cor-high-derivative-signed} on the actual central
section. Its complete derivative bound holds on this entire interval,
including \(h_L=0\) and the enlarged source chord through that orbit.
Take the maximum of the compact and high-height derivative constants,
then choose \(m\) so that \(\Delta'>1/2\) on both intervals.
The annular curvature comes from Lemma~\ref{bd:lem-annulus}, on its
open overlap with the source interval constructed by the two fixed
comparisons. All assertions concern restrictions of one physical
map and every input in those overlaps. Boundary coefficients denote
normalized one-sided limits; no orbit at a collapsed raw section is
asserted.
\end{proof}
